\documentclass[12pt]{report}
\usepackage[left=1in, right=0.5in, top=0.5in, bottom=0.5in]{geometry}
\usepackage{listings}
\usepackage{xcolor}
\usepackage{fontspec}
\usepackage{caption}
\usepackage{graphicx}
\usepackage{hyperref}

\setmainfont{Sarabun}

% Colors
\definecolor{codegreen}{rgb}{0,0.6,0}
\definecolor{codegray}{rgb}{0.5,0.5,0.5}
\definecolor{codepurple}{rgb}{0.58,0,0.82}
\definecolor{backcolour}{rgb}{0.95,0.95,0.92}

\lstdefinestyle{mystyle}{
    backgroundcolor=\color{backcolour},   
    commentstyle=\color{codegreen},
    keywordstyle=\color{magenta},
    numberstyle=\tiny\color{codegray},
    stringstyle=\color{codepurple},
    basicstyle=\ttfamily\footnotesize,
    breakatwhitespace=false,         
    breaklines=true,                 
    captionpos=b,                    
    keepspaces=true,                 
    numbers=left,                    
    numbersep=5pt,                  
    showspaces=false,                
    showstringspaces=false,
    showtabs=false,                  
    tabsize=2
}

\lstset{style=mystyle}

% JavaScript Definition
\lstdefinelanguage{JavaScript}{
  keywords={break, case, catch, class, const, continue, debugger, default, delete, do, else, enum, export, extends, false, finally, for, function, if, import, in, instanceof, new, null, return, super, switch, this, throw, true, try, typeof, var, void, while, with, yield, async, await},
  keywordstyle=\color{blue}\bfseries,
  ndkeywords={class, export, boolean, throw, implements, import, this},
  ndkeywordstyle=\color{darkgray}\bfseries,
  identifierstyle=\color{black},
  sensitive=false,
  comment=[l]{//},
  morecomment=[s]{/*}{*/},
  commentstyle=\color{purple}\ttfamily,
  stringstyle=\color{red}\ttfamily,
  morestring=[b]',
  morestring=[b]"
}

\title{Term Project: GymBro (Queue Booking System)}
\author{Full Stack Development}
\date{\today}

\begin{document}

% Title Page
\begin{titlepage}
    \centering
    \vspace*{1cm}
    \Huge
    \textbf{Term Project: GymBro}
    
    \vspace{0.5cm}
    \LARGE
    Queue Booking System
    
    \vspace{1.5cm}
    \textbf{Course:} Full Stack Development
    
    \vfill
    \Large
    \today
    
\end{titlepage}

\tableofcontents
\newpage

\chapter{บทนำ (Introduction)}

\section{ภาพรวมของโครงการ (Project Overview)}
GymBro คือระบบจองคิวและบริหารจัดการยิมที่พัฒนาขึ้นเพื่ออำนวยความสะดวกให้กับทั้งผู้ดูแลระบบ (Admin) และสมาชิก (User) โดยระบบรองรับฟังก์ชันการทำงานหลักดังนี้:
\begin{itemize}
    \item \textbf{ระบบสมาชิก (Membership System)}: สมัครสมาชิก, ล็อกอิน, จัดการข้อมูลส่วนตัว
    \item \textbf{ระบบจองคิว (Booking System)}: จองเซสชันกับเทรนเนอร์และอุปกรณ์, ตรวจสอบตารางเวลาว่าง
    \item \textbf{ระบบจัดการหลังบ้าน (Admin Dashboard)}: จัดการข้อมูลลูกค้า, เทรนเนอร์, อุปกรณ์ และดูรายงานสรุป
\end{itemize}

\section{โครงสร้างของโปรเจกต์ (Project Structure)}
โปรเจกต์นี้ถูกออกแบบโดยแยกส่วนการทำงานระหว่าง Backend และ Frontend อย่างชัดเจน (Separation of Concerns) ดังนี้:

\begin{itemize}
    \item \textbf{Backend}: พัฒนาด้วย Node.js และ Express.js เชื่อมต่อกับฐานข้อมูล PostgreSQL ผ่าน Sequelize ORM
    \item \textbf{Frontend}: พัฒนาด้วย HTML, Tailwind CSS, และ Client-side JavaScript โดยมี Node.js Server ทำหน้าที่ให้บริการไฟล์ Static
\end{itemize}

\chapter{การออกแบบฐานข้อมูล (Database Design)}

\section{แผนภาพความสัมพันธ์ (Database Schema)}
ฐานข้อมูลถูกออกแบบให้มีความสัมพันธ์แบบ Relational Database โดยมีตารางหลักคือ Users (Customers), Trainers, Equipments, และ Sessions

\subsection{SQL Schema (backend/schema.sql)}
\begin{lstlisting}[language=SQL, caption={โครงสร้างฐานข้อมูลใน backend/schema.sql}, label={code:schema-sql}]
-- Drop tables if exists
DROP TABLE IF EXISTS "Sessions";
DROP TABLE IF EXISTS "Equipments";
DROP TABLE IF EXISTS "Trainers";
DROP TABLE IF EXISTS "Customers";

-- 1. Customers Table
CREATE TABLE "Customers" (
    id SERIAL PRIMARY KEY,
    "fullName" VARCHAR(255) NOT NULL,
    email VARCHAR(255) UNIQUE NOT NULL,
    phone VARCHAR(20),
    "memberType" VARCHAR(50) DEFAULT 'Standard', -- Standard, Premium, VIP
    "memberLevel" VARCHAR(50) DEFAULT 'Basic',   -- Basic, Silver, Gold
    "startDate" TIMESTAMP WITH TIME ZONE DEFAULT CURRENT_TIMESTAMP,
    "endDate" TIMESTAMP WITH TIME ZONE,
    "createdAt" TIMESTAMP WITH TIME ZONE DEFAULT CURRENT_TIMESTAMP,
    "updatedAt" TIMESTAMP WITH TIME ZONE DEFAULT CURRENT_TIMESTAMP
);

-- 2. Trainers Table
CREATE TABLE "Trainers" (
    id SERIAL PRIMARY KEY,
    name VARCHAR(255) NOT NULL,
    specialty VARCHAR(100), -- Yoga, Cardio, Bodybuilding
    level VARCHAR(50) DEFAULT 'Junior', -- Junior, Senior, Master
    phone VARCHAR(20),
    "createdAt" TIMESTAMP WITH TIME ZONE DEFAULT CURRENT_TIMESTAMP,
    "updatedAt" TIMESTAMP WITH TIME ZONE DEFAULT CURRENT_TIMESTAMP
);

-- 3. Equipments Table
CREATE TABLE "Equipments" (
    id SERIAL PRIMARY KEY,
    name VARCHAR(255) NOT NULL,
    category VARCHAR(100), -- Treadmill, Dumbbell, Machine
    status VARCHAR(50) DEFAULT 'Available', -- Available, Maintenance, In Use
    "createdAt" TIMESTAMP WITH TIME ZONE DEFAULT CURRENT_TIMESTAMP,
    "updatedAt" TIMESTAMP WITH TIME ZONE DEFAULT CURRENT_TIMESTAMP
);

-- 4. Sessions Table (Booking)
CREATE TABLE "Sessions" (
    id SERIAL PRIMARY KEY,
    "sessionDate" TIMESTAMP WITH TIME ZONE NOT NULL,
    duration INTEGER DEFAULT 60, -- Minutes
    status VARCHAR(50) DEFAULT 'Scheduled', -- Scheduled, Completed, Cancelled
    "customerId" INTEGER REFERENCES "Customers"(id) ON DELETE CASCADE,
    "trainerId" INTEGER REFERENCES "Trainers"(id) ON DELETE SET NULL,
    "equipmentId" INTEGER REFERENCES "Equipments"(id) ON DELETE SET NULL,
    "createdAt" TIMESTAMP WITH TIME ZONE DEFAULT CURRENT_TIMESTAMP,
    "updatedAt" TIMESTAMP WITH TIME ZONE DEFAULT CURRENT_TIMESTAMP
);
\end{lstlisting}

\section{การเชื่อมต่อฐานข้อมูล (ORM Implementation)}
ใช้ Sequelize ในการจัดการ Model และ Association ระหว่างตาราง

\subsection{Sequelize Setup (backend/db.js)}
\begin{lstlisting}[language=JavaScript, caption={การตั้งค่า Sequelize และ Models ใน backend/db.js}, label={code:backend-db}]
const { Sequelize, DataTypes } = require('sequelize');

// Connection Setup
const sequelize = new Sequelize('gymbro_db', 'postgres', 'postgres', {
    host: 'localhost',
    dialect: 'postgres',
    logging: false
});

// 1. Customer Model
const Customer = sequelize.define('Customer', {
    fullName: { type: DataTypes.STRING, allowNull: false },
    email: { type: DataTypes.STRING, unique: true, allowNull: false },
    phone: { type: DataTypes.STRING },
    memberType: { type: DataTypes.STRING, defaultValue: 'Standard' },
    memberLevel: { type: DataTypes.STRING, defaultValue: 'Basic' },
    startDate: { type: DataTypes.DATE, defaultValue: DataTypes.NOW },
    endDate: { type: DataTypes.DATE }
});

// 2. Trainer Model
const Trainer = sequelize.define('Trainer', {
    name: { type: DataTypes.STRING, allowNull: false },
    specialty: { type: DataTypes.STRING },
    level: { type: DataTypes.STRING, defaultValue: 'Junior' },
    phone: { type: DataTypes.STRING }
});

// 3. Equipment Model
const Equipment = sequelize.define('Equipment', {
    name: { type: DataTypes.STRING, allowNull: false },
    category: { type: DataTypes.STRING },
    status: { type: DataTypes.STRING, defaultValue: 'Available' }
});

// 4. Session Model
const Session = sequelize.define('Session', {
    sessionDate: { type: DataTypes.DATE, allowNull: false },
    duration: { type: DataTypes.INTEGER, defaultValue: 60 },
    status: { type: DataTypes.STRING, defaultValue: 'Scheduled' }
});

// Associations
Customer.hasMany(Session, { foreignKey: 'customerId' });
Session.belongsTo(Customer, { foreignKey: 'customerId' });

Trainer.hasMany(Session, { foreignKey: 'trainerId' });
Session.belongsTo(Trainer, { foreignKey: 'trainerId' });

Equipment.hasMany(Session, { foreignKey: 'equipmentId' });
Session.belongsTo(Equipment, { foreignKey: 'equipmentId' });

module.exports = { sequelize, Customer, Trainer, Equipment, Session };
\end{lstlisting}


\chapter{การทำงานฝั่ง Server (Backend API)}

\section{การตั้งค่า Server และ API Routes (backend/server.js)}
ไฟล์ \texttt{backend/server.js} เป็นจุดเริ่มต้นของ Backend Application ซึ่งทำหน้าที่เชื่อมต่อฐานข้อมูลและให้บริการ REST API สำหรับ Frontend

\begin{lstlisting}[language=JavaScript, caption={การตั้งค่า Express Server และ API Routes}, label={code:backend-server}]
const express = require('express');
const cors = require('cors');
const { sequelize, Customer, Trainer, Equipment, Session } = require('./db');
const app = express();
const port = 3000;

app.use(cors());
app.use(express.json());

// Sync Database
sequelize.sync({ force: false }).then(() => {
    console.log('Database synced');
});

// ==========================================
// API Routes
// ==========================================

// 1. Customers API
app.get('/api/customers', async (req, res) => {
    const customers = await Customer.findAll();
    res.json(customers);
});

app.post('/api/customers', async (req, res) => {
    try {
        const customer = await Customer.create(req.body);
        res.status(201).json(customer);
    } catch (err) {
        res.status(400).json({ error: err.message });
    }
});

// 2. Trainers API
app.get('/api/trainers', async (req, res) => {
    const trainers = await Trainer.findAll();
    res.json(trainers);
});

// 3. Equipments API
app.get('/api/equipments', async (req, res) => {
    const equipments = await Equipment.findAll();
    res.json(equipments);
});

// 4. Sessions API (Booking Logic)
app.get('/api/sessions', async (req, res) => {
    const { date, today } = req.query;
    const where = {};
    
    if (today) {
        const start = new Date();
        start.setHours(0,0,0,0);
        const end = new Date();
        end.setHours(23,59,59,999);
        where.sessionDate = { [Sequelize.Op.between]: [start, end] };
    } else if (date) {
        const start = new Date(date);
        const end = new Date(date);
        end.setDate(end.getDate() + 1);
        where.sessionDate = { [Sequelize.Op.between]: [start, end] };
    }

    const sessions = await Session.findAll({
        where,
        include: [Customer, Trainer, Equipment],
        order: [['sessionDate', 'ASC']]
    });
    res.json(sessions);
});

app.post('/api/sessions', async (req, res) => {
    // Booking Logic: Check availability
    // (Implementation omitted for brevity)
    const session = await Session.create(req.body);
    res.status(201).json(session);
});

// Start Server
app.listen(port, () => {
    console.log(`Backend server running on http://localhost:${port}`);
});
\end{lstlisting}

\chapter{ส่วนติดต่อผู้ใช้ (Frontend Application)}

\section{Frontend Server (frontend/server.js)}
ทำหน้าที่ให้บริการไฟล์ Static และ Routing เบื้องต้นสำหรับหน้าเว็บต่างๆ

\begin{lstlisting}[language=JavaScript, caption={Frontend Server Setup}, label={code:frontend-server}]
const express = require('express');
const path = require('path');
const app = express();
const port = 5500;

app.set('view engine', 'ejs');
app.set('views', path.join(__dirname, 'views'));
app.use(express.static(path.join(__dirname, 'public')));
app.use('/js', express.static(path.join(__dirname, 'js')));

// Routes
app.get('/', (req, res) => res.redirect('/login'));
app.get('/login', (req, res) => res.render('login'));
app.get('/register', (req, res) => res.render('register'));

// Admin Routes
app.get('/admin/dashboard', (req, res) => res.render('index'));
app.get('/admin/customers', (req, res) => res.render('customers'));
// ... other admin routes

// User Routes
app.get('/user/dashboard', (req, res) => res.render('user/dashboard'));
app.get('/user/profile', (req, res) => res.render('user/profile'));

app.listen(port, () => {
    console.log(`Frontend server running on http://localhost:${port}`);
});
\end{lstlisting}

\section{การจัดการสิทธิ์และ Configuration (Core JS)}

\subsection{Configuration (frontend/js/config.js)}
\begin{lstlisting}[language=JavaScript, caption={Global Configuration}, label={code:frontend-config}]
window.API = "http://localhost:3000/api";
\end{lstlisting}

\subsection{Authentication Logic (frontend/js/auth.js)}
จัดการตรวจสอบสถานะการล็อกอินและ Role-based Access Control

\begin{lstlisting}[language=JavaScript, caption={Authentication Guard}, label={code:frontend-auth}]
const isLoginPage = window.location.pathname === '/login' || window.location.pathname === '/';

if (!isLoginPage) {
    const session = JSON.parse(localStorage.getItem('gymbro_user'));
    if (!session || !session.user) {
        window.location.href = '/login';
    } else {
        // Role Guard
        const path = window.location.pathname;
        if (path.startsWith('/admin') && session.role !== 'admin') {
            window.location.href = '/user/dashboard';
        }
    }
}

window.logout = function() {
    localStorage.removeItem('gymbro_user');
    window.location.href = '/login';
};
\end{lstlisting}


\chapter{ส่วนแสดงผลและส่วนติดต่อผู้ใช้ (User Interface Views)}

\section{ระบบยืนยันตัวตน (Authentication)}

\subsection{หน้าเข้าสู่ระบบ (frontend/views/login.ejs)}
\begin{lstlisting}[language=HTML, caption={Login Page Structure}, label={code:views-login}]
<body class="bg-primary flex items-center justify-center min-h-screen p-4">
  <div class="w-full max-w-md bg-contrast shadow-[8px_8px_0px_0px_#629FAD] border-4 border-secondary p-8 relative">
    <div class="text-center mb-8">
      <h1 class="text-4xl font-extrabold text-primary uppercase italic tracking-tighter mb-2">Gym<span class="text-secondary">Bro</span></h1>
      <p class="text-secondary font-bold text-sm uppercase tracking-widest">Queue Booking System</p>
    </div>

    <form id="login-form" class="space-y-6">
      <!-- Email & Phone Inputs -->
      <button type="submit" class="w-full bg-primary text-contrast font-black uppercase py-4 hover:bg-secondary transition-all">
        Sign In
      </button>
    </form>
  </div>
  <script src="/js/auth.js"></script>
</body>
\end{lstlisting}

\subsection{หน้าลงทะเบียน (frontend/views/register.ejs)}
\begin{lstlisting}[language=HTML, caption={Register Page Structure}, label={code:views-register}]
<form id="register-form" class="space-y-4">
  <input type="text" id="fullName" placeholder="Full Name" required>
  <input type="email" id="email" placeholder="Email Address" required>
  <input type="tel" id="phone" placeholder="Phone Number" required>
  
  <button type="submit" class="w-full bg-primary text-contrast font-black uppercase py-3">
    Create Account
  </button>
</form>

<script>
  // Inline Script for Registration
  const form = document.getElementById('register-form');
  form.addEventListener('submit', async (e) => {
    e.preventDefault();
    // Fetch API Logic...
  });
</script>
\end{lstlisting}

\section{ระบบจัดการหลังบ้าน (Admin Dashboard)}

\subsection{หน้าหลัก (frontend/views/index.ejs)}
แสดงภาพรวมสถิติและตารางงานวันนี้

\begin{lstlisting}[language=HTML, caption={Admin Dashboard Layout}, label={code:views-index}]
<!-- Sidebar -->
<aside class="w-64 bg-primary text-contrast hidden md:flex flex-col">
  <nav class="flex-1 overflow-y-auto py-6">
    <!-- Menu Items -->
  </nav>
</aside>

<!-- Main Content -->
<main class="flex-1 flex flex-col h-screen overflow-hidden relative">
  <!-- Stats Grid -->
  <div class="grid grid-cols-1 md:grid-cols-4 gap-6 mb-8">
    <div class="bg-white border-2 border-primary p-6 shadow-[4px_4px_0px_0px_#0C2C55]">
      <h3 class="text-3xl font-black text-primary" id="count-customers">0</h3>
      <p class="text-sm font-bold text-gray-500 uppercase">Members</p>
    </div>
  </div>

  <!-- Today's Schedule Table -->
  <div class="bg-white border-2 border-primary shadow-[8px_8px_0px_0px_#296374]">
    <table>
      <tbody id="today-sessions"></tbody>
    </table>
  </div>
</main>
\end{lstlisting}


\section{การจัดการทรัพยากร (Resource Management)}

\subsection{การจัดการสมาชิก (frontend/views/customers.ejs)}
ระบบจัดการสมาชิก รองรับการเพิ่ม, ลบ, แก้ไข และแสดงรายชื่อสมาชิกทั้งหมด

\begin{lstlisting}[language=HTML, caption={Customer Management Interface}, label={code:views-customers}]
<!-- Add Customer Form -->
<form id="customer-form" class="grid grid-cols-1 md:grid-cols-4 gap-4 mb-8">
  <input id="name" placeholder="Full Name" required>
  <input id="email" type="email" placeholder="Email" required>
  <select id="memberType">
    <option value="Standard">Standard</option>
    <option value="Premium">Premium</option>
  </select>
  <button type="submit" class="bg-primary text-contrast font-bold uppercase">Save</button>
</form>

<!-- Customer List Table -->
<table class="w-full text-left border-collapse">
  <thead>
    <tr>
      <th>ID</th><th>Name</th><th>Type</th><th>Actions</th>
    </tr>
  </thead>
  <tbody id="customers-table">
    <!-- Populated by JS -->
  </tbody>
</table>
<script src="/js/customers.js"></script>
\end{lstlisting}

\subsection{การจัดการเทรนเนอร์ (frontend/views/trainers.ejs)}
\begin{lstlisting}[language=HTML, caption={Trainer Management Interface}, label={code:views-trainers}]
<!-- Add Trainer Form -->
<form id="trainer-form" class="grid grid-cols-1 md:grid-cols-5 gap-4 mb-8">
  <input id="name" placeholder="Name" required>
  <input id="specialty" placeholder="Specialty">
  <select id="trainerLevel">
    <option value="Junior">Junior</option>
    <option value="Senior">Senior</option>
  </select>
  <button type="submit">Save Trainer</button>
</form>

<!-- Trainer List Table -->
<table class="w-full">
  <tbody id="trainers-table"></tbody>
</table>
<script src="/js/trainers.js"></script>
\end{lstlisting}

\subsection{การจัดการอุปกรณ์ (frontend/views/equipments.ejs)}
\begin{lstlisting}[language=HTML, caption={Equipment Management Interface}, label={code:views-equipments}]
<!-- Add Equipment Form -->
<form id="equipment-form" class="grid grid-cols-1 md:grid-cols-4 gap-4 mb-8">
  <input id="name" placeholder="Name" required>
  <input id="category" placeholder="Category">
  <select id="status">
    <option value="Available">Available</option>
    <option value="Maintenance">Maintenance</option>
  </select>
  <button type="submit">Save Equipment</button>
</form>

<!-- Equipment List Table -->
<table class="w-full">
  <tbody id="equipments-table"></tbody>
</table>
<script src="/js/equipments.js"></script>
\end{lstlisting}

\section{การจัดการเซสชัน (Session Management)}

\subsection{หน้าจอจัดการการจอง (frontend/views/sessions.ejs)}
สำหรับผู้ดูแลระบบในการสร้างการจองและดูตารางทั้งหมด

\begin{lstlisting}[language=HTML, caption={Session Management Interface}, label={code:views-sessions}]
<!-- Booking Form -->
<form id="session-form" class="grid grid-cols-1 md:grid-cols-5 gap-4">
  <select id="customer" required><option>Select Customer</option></select>
  <select id="trainer" required><option>Select Trainer</option></select>
  <select id="equipment" required><option>Select Equipment</option></select>
  <input id="scheduled_at" type="datetime-local" required>
  <button type="submit">Book Now</button>
</form>

<!-- Sessions List Table -->
<table class="w-full">
  <thead>
    <tr><th>Customer</th><th>Trainer</th><th>Date</th><th>Actions</th></tr>
  </thead>
  <tbody id="sessions-table"></tbody>
</table>
<script src="/js/sessions.js"></script>
\end{lstlisting}


\section{รายงานและการตรวจสอบ (Reports \& Logs)}

\subsection{หน้าแสดงรายงาน (frontend/views/logs.ejs)}
หน้าจอแสดงบันทึกการใช้งานและรายงานสรุปสำหรับผู้ดูแลระบบ

\begin{lstlisting}[language=HTML, caption={Logs and Reports Interface}, label={code:views-logs}]
<!-- Tab Controls -->
<div class="flex">
  <button onclick="switchTab('logs')" class="tab-active">System Logs</button>
  <button onclick="switchTab('customer')" class="tab-inactive">Customer Report</button>
</div>

<!-- Log Tables (Hidden by default) -->
<div id="section-logs" class="bg-white p-4">
  <table class="w-full">
    <thead><tr><th>Time</th><th>Action</th><th>Details</th></tr></thead>
    <tbody id="logs-table"></tbody>
  </table>
</div>
<script src="/js/logs.js"></script>
\end{lstlisting}

\chapter{ส่วนสมาชิก (User Portal)}

\section{แดชบอร์ดสมาชิก (frontend/views/user/dashboard.ejs)}
หน้าหลักสำหรับสมาชิกเพื่อดูตารางฝึกซ้อมและจองเซสชันด้วยตนเอง

\begin{lstlisting}[language=HTML, caption={User Dashboard Interface}, label={code:views-user-dashboard}]
<!-- Booking Modal -->
<div id="booking-modal" class="hidden fixed inset-0 bg-black/50 z-50">
  <div class="bg-white p-6 max-w-md mx-auto mt-20">
    <h3 class="text-xl font-bold">Book New Session</h3>
    <form id="booking-form">
      <input type="date" id="book-date" required>
      <select id="book-trainer"></select>
      <button type="submit">Confirm</button>
    </form>
  </div>
</div>

<!-- Schedule Table -->
<div class="bg-white border-2 border-primary">
  <h3 class="bg-secondary text-white p-4">All Bookings</h3>
  <table class="w-full">
    <tbody id="schedule-table-body"></tbody>
  </table>
</div>
<script src="/js/user/dashboard.js"></script>
\end{lstlisting}

\section{โปรไฟล์สมาชิก (frontend/views/user/profile.ejs)}
หน้าแสดงข้อมูลส่วนตัว สถานะสมาชิก และประวัติการใช้งาน

\begin{lstlisting}[language=HTML, caption={User Profile Interface}, label={code:views-user-profile}]
<div class="flex flex-col md:flex-row gap-8">
  <!-- Avatar Card -->
  <div class="flex-shrink-0 flex flex-col items-center">
    <div class="w-32 h-32 bg-secondary text-white flex items-center justify-center">
      <i data-lucide="user"></i>
    </div>
    <span class="bg-primary text-accent px-3 py-1">Active Member</span>
  </div>

  <!-- Details -->
  <div class="flex-1 space-y-6">
    <div>
      <label>Full Name</label>
      <h3 id="profile-name" class="text-2xl font-black">Loading...</h3>
    </div>
    <div>
      <label>Membership Type</label>
      <span id="profile-type" class="bg-contrast px-4 py-2 border-2 border-primary">Standard</span>
    </div>
  </div>
</div>
<script src="/js/user/profile.js"></script>
\end{lstlisting}

\end{document}
